\subsection{Groundwater flow calibration with a baseflow constraint}
\label{sec:calibration}

The calibrated flow field is simultaneously head- and flux-consistent: the deployed posterior
realization reproduces $5{,}383$ static water-table heads at a Nash--Sutcliffe efficiency of $0.91$
(RMSE $5.6$~m) while also reproducing the basin's gaining baseflow, a joint property a head-only fit
cannot guarantee. This is achieved by calibrating the groundwater flow model to both water-table
elevations and net baseflow with the PEST\texttt{++} iterative ensemble smoother
\citep{white2020pestpp}.

The adjustable parameters (Table~\ref{tab:calparams}) are a smooth hydraulic-conductivity field of
$243$ pilot-point log-multipliers interpolated to the grid, supplemented by global multipliers on
vertical anisotropy, recharge, drain conductance, the watershed-divide boundary conductance, well
pumping, and streambed conductance (the SFR analogue of river conductance). Observations are
$5{,}383$ static water-table heads --- the per-cell median of the \emph{raw} measured static water
levels in the well database ($\text{head} = \text{top} - \text{SWL}$), not values drawn from the
kriged water-table surface --- plus a single net-baseflow observation. The targets are therefore
independent measurements rather than a smooth product, with one noted partial non-independence: the
general-head boundary is set from the kriged surface derived from the same database. We emphasise the
RMSE ($5.6$~m) over the Nash--Sutcliffe efficiency ($0.91$), since the latter is partly carried by
the large topographic range of heads.

The baseflow constraint is essential. A head-only inversion is non-unique and can fit heads while
producing an unphysical (net-losing) stream network, so the basin groundwater-to-stream exchange
(from the SFR budget) is added as a strongly weighted target equal to the observed gaining baseflow
($5.56$~m\textsuperscript{3}~s\textsuperscript{-1}; USGS~04118500 mean flow
$7.51$~m\textsuperscript{3}~s\textsuperscript{-1}, baseflow index $0.74$). Two domain controls keep
the constrained problem well posed: raw well \emph{pump capacity} overstates long-term withdrawal in
this heavily developed basin and is scaled by a calibrated multiplier, and the watershed perimeter is
treated as a near-no-flow divide (a topographic divide is a groundwater flow divide; a leaky boundary
exports baseflow out the sides). The ensemble converges in ten minutes on a single workstation; the
posterior realization closest to the observed baseflow gives the head Nash--Sutcliffe efficiency of
$0.91$ (RMSE $5.6$~m over $5{,}383$ heads) and a gaining baseflow of
$+5.46$~m\textsuperscript{3}~s\textsuperscript{-1}. This flow field is the basis for the groundwater
transport.
